% ============================================================
%  Section: Visualization
% ============================================================

The results of the analysis are presented through two complementary tools:
a \textbf{Tableau Public dashboard} for static, publication-quality charts,
and a \textbf{Streamlit web application} for interactive querying.

\subsection{Tableau Dashboard}

% Insert screenshots once your Tableau workbook is complete.
\textit{[Add Tableau screenshots and description here]}

\subsection{Streamlit Application}

% Insert screenshots and describe the app's pages/features.
\textit{[Add Streamlit screenshots and description here]}

% ============================================================
%  Section: Conclusion
% ============================================================
% (This file is input twice in main.tex — once for visualization,
%  once for conclusion. Split into two files when you fill them in.)

\subsection{Summary}

\textit{[Write a 1-2 paragraph summary of what you built and what you found.]}

\subsection{Challenges and Lessons Learned}

\textit{[Describe the hardest parts: data quality issues, SQL problems you debugged,
anything that didn't go as planned and how you resolved it.]}

\subsection{Future Work}

If this project were extended, the following enhancements would add significant value:

\begin{itemize}
    \item Integration of real-time GTFS delay feeds to correlate ridership drops
          with service disruptions.
    \item Weather data enrichment to test the hypothesis that rain increases
          ridership on certain lines.
    \item A PostgreSQL migration to support multi-user access and larger datasets.
    \item Scheduled ingestion via a cloud function (AWS Lambda or GitHub Actions)
          so the database updates automatically every Monday when MTA publishes
          the new weekly file.
\end{itemize}
